\section*{Chapitre \ref{ch:b3:group-theory-applied} --- La théorie des groupes en action}

\begin{solution}{exo:b3:group-theory-applied:1}
$n_{A_1} = \frac14(5 - 1 + 3 + 1) = 2$, $n_{A_2} = \frac14(5 - 1 - 3 - 1) = 0$,
$n_{B_1} = \frac14(5 + 1 + 3 - 1) = 2$, $n_{B_2} = \frac14(5 + 1 - 3 + 1) = 1$~:
$\Gamma = 2A_1 + 2B_1 + B_2$, de dimension 5.
\end{solution}

\begin{solution}{exo:b3:group-theory-applied:2}
Même décompte des atomes laissés en place que pour l'eau~: $\Gamma_{\mathrm{vib}} = 2A_1 + B_1$.
$A_1$ ($z$) et
$B_1$ ($x$) sont \omterm{def:b3:group-theory-applied:activity}{actifs en IR}, et tous deux sont aussi \omterm{def:b3:group-theory-applied:activity}{actifs en Raman}
($x^2$, $xz$)~: trois bandes dans chaque spectre.
\end{solution}

\begin{solution}{exo:b3:group-theory-applied:3}
$B_1\otimes B_2$~: \omterm{def:b3:point-groups:representation}{caractères} $(1, 1, -1, -1)$ $= A_2$. $A_2\otimes B_1$~: $(1, -1, -1,
1) = B_2$. Dans $C_{3v}$, $E\otimes E$ a pour \omterm{def:b3:point-groups:representation}{caractères} $(4, 1, 0)$~;
réduction~: $A_1 + A_2 +
E$.
\end{solution}

\begin{solution}{exo:b3:group-theory-applied:4}
Infrarouge~: les deux modes $t_2$, 3019 et \qty{1306}{cm^{-1}}. Raman~:
les quatre ($a_1$,
$e$, $t_2$ se transforment tous comme des fonctions quadratiques). Pas de
\omterm{def:b3:point-groups:inversion}{centre d'inversion}, donc des coïncidences sont possibles.
\end{solution}

\begin{solution}{exo:b3:group-theory-applied:5}
Les atomes laissés en place $(4, 1, 2, 4, 1, 2)$ multipliés par $\chi_{xyz} = (3, 0, -1, 1, -2, 1)$ donnent
$\chi_{3N} = (12, 0, -2, 4, -2, 2)$, qui se réduit en $A_1' + A_2' + 3E' + 2A_2'' +
E''$. Retirer $E' + A_2''$ (translations) et $A_2' + E''$ (rotations)~: $A_1' + 2E'
+ A_2''$. IR~: $2E' + A_2''$, trois bandes~; Raman~: $A_1' + 2E'$, trois
bandes. $A_2''$
(déformation hors du plan) n'est actif qu'en IR, $A_1'$ (élongation
symétrique) qu'en Raman, les modes $E'$ apparaissent dans les deux.
\end{solution}

\begin{solution}{exo:b3:group-theory-applied:6}
IR~: les deux modes $T_{1u}$, deux bandes. Raman~: $A_{1g}$, $E_g$,
$T_{2g}$, trois bandes.
$T_{2u}$ est silencieux. \ce{SF6} possède un \omterm{def:b3:point-groups:inversion}{centre d'inversion}~:
exclusion mutuelle.
\end{solution}

\begin{solution}{exo:b3:group-theory-applied:7}
Par $E$, $C_3$, $C_3^2$, la fonction $h_2$ devient $h_2$, $h_3$, $h_1$
(\omterm{def:b3:point-groups:representation}{caractères} 2,
$-1$, $-1$)~; les réflexions contribuent pour 0. $\hat P_Eh_2 \propto 2h_2 - h_3 - h_1$.
Son recouvrement avec $2h_1 - h_2 - h_3$ (recouvrements atomiques
négligés) vaut $-3$, et cette dernière a pour norme $\sqrt6$~; en
retranchant la projection, $(2h_2 - h_1 - h_3) +
\frac12(2h_1 - h_2 - h_3) = \frac32(h_2 - h_3)$.
\end{solution}

\begin{solution}{exo:b3:group-theory-applied:8}
fac ($C_{3v}$, \omterm{def:b3:point-groups:representation}{caractères} $3, 0, 1$)~: $A_1 + E$, deux bandes IR. mer
($C_{2v}$, \omterm{def:b3:point-groups:representation}{caractères} $3, 1, 3, 1$)~: $2A_1 + B_1$, trois bandes.
\ce{M(CO)6} ($O_h$)~: $A_{1g} + E_g
+ T_{1u}$, une bande IR ($T_{1u}$).
\end{solution}

\begin{solution}{exo:b3:group-theory-applied:9}
\omterm{def:b3:point-groups:representation}{Caractères} $6, 0, 0, 2, 2, 0, 0, 0, 4, 2$~: seuls $E$, les $C_4$ et les
$C_2$ portés par les axes (deux ligands laissés en place), les trois
$\sigma_h$ (quatre) et les six $\sigma_d$
(deux) laissent des ligands en place. La réduction donne
$A_{1g} + E_g + T_{1u}$. Les orbitales du métal $s$
($a_{1g}$), $d_{z^2}$ et $d_{x^2-y^2}$ ($e_g$), et $p_x$, $p_y$, $p_z$
($t_{1u}$) se combinent~;
$t_{2g}$ reste non liant.
\end{solution}

\begin{solution}{exo:b3:group-theory-applied:10}
$n_i = \frac1h(1\cdot h\cdot\chi_i(E)) = d_i$. La dimension de la
représentation régulière vaut $h$, et aussi $\sum_in_id_i = \sum_id_i^2$.
\end{solution}

\begin{solution}{exo:b3:group-theory-applied:11}
Dans $D_{2h}$ (molécule selon $z$), atomes laissés en place $(4, 4, 0, 0, 0, 0, 4, 4)$, $\chi_{3N}
= (12, -4, 0, 0, 0, 0, 4, 4)$, qui se réduit en $2A_g + 2B_{2g} + 2B_{3g} + 2B_{1u} + 2B_{2u}
+ 2B_{3u}$. Retirer les translations $B_{1u} + B_{2u} + B_{3u}$ et les
deux rotations
$B_{2g} + B_{3g}$~: $2A_g + B_{2g} + B_{3g} + B_{1u} + B_{2u} + B_{3u}$,
sept modes.
Dans $D_{\infty h}$~: $2\Sigma_g^+ + \Pi_g + \Sigma_u^+ + \Pi_u$. IR~:
$\Sigma_u^+$ et $\Pi_u$
(deux bandes)~; Raman~: $2\Sigma_g^+$ et $\Pi_g$ (trois bandes)~; aucune
coïncidence.
\end{solution}

\begin{solution}{exo:b3:group-theory-applied:12}
$1b_2 \to 4a_1$~: $B_2$, $y$, permise, perpendiculaire au plan. $3a_1 \to 4a_1$~:
$A_1$, $z$, permise, selon l'axe $C_2$. $1b_2 \to 2b_1$~: $A_2$, interdite.
$1b_1 \to 4a_1$~: $B_1$, $x$, permise, dans le plan, perpendiculairement
à l'axe.
\end{solution}

\begin{solution}{pb:b3:group-theory-applied:1}
\textbf{1.} cis~: les deux CO sur des sommets adjacents~; trans~:
opposés~; fac~: trois CO sur une face~; mer~: trois CO sur un méridien
(deux en trans, un entre les deux).
\textbf{2.} $C_{2v}$.
\textbf{3.} $D_{4h}$.
\textbf{4.} $C_{3v}$.
\textbf{5.} $C_{2v}$.
\textbf{6.} Seul le trans-\ce{ML4(CO)2}.
\textbf{7.} Une opération qui envoie un CO sur un autre place son vecteur
hors de la diagonale (contribution 0)~; une opération qui laisse un CO en
place envoie son vecteur de liaison sur lui-même (+1)~; aucune ne
retourne un vecteur C--O.
\textbf{8.} \omterm{def:b3:point-groups:representation}{Caractères} $(2, 0, 2, 0)$, avec $\sigma_v(xz)$ le plan des
deux CO~: $A_1 +
B_1$.
\textbf{9.} \omterm{def:b3:point-groups:representation}{Caractères} 2 pour $E$, $C_4$, $C_2$, $\sigma_v$, $\sigma_d$,
0 ailleurs~:
$A_{1g} + A_{2u}$.
\textbf{10.} $(3, 0, 1)$~: $A_1 + E$.
\textbf{11.} $(3, 1, 3, 1)$~: $2A_1 + B_1$.
\textbf{12.} 2, 2, 3, 3~: une dimension par CO.
\textbf{13.} cis 2, trans 1, fac 2, mer 3.
\textbf{14.} cis 2, trans 1, fac 2, mer 3.
\textbf{15.} trans~: sa bande IR ($A_{2u}$) et sa bande Raman ($A_{1g}$)
sont des modes différents.
\textbf{16.} $A_{2u}$ change de signe par $i$~: un CO s'allonge pendant
que l'autre se raccourcit, en opposition de phase.
\textbf{17.} $E$ est une paire dégénérée~: deux modes de même fréquence
donnent une seule bande, plus la bande $A_1$.
\textbf{18.} Un isomère trans-disubstitué, centrosymétrique.
\textbf{19.} Trois bandes IR~: mer (fac en donnerait deux).
\textbf{20.} Oui~: mer, n'étant pas centrosymétrique, a trois modes CO
\omterm{def:b3:group-theory-applied:activity}{actifs en Raman} aux mêmes fréquences~; fac en montrerait deux.
\textbf{21.} $A_1$~: l'élongation en phase est totalement symétrique.
\textbf{22.} Un mélange fac/mer montrerait jusqu'à cinq bandes IR, les
bandes de fac à d'autres positions~; trois bandes nettes correspondent à
un seul isomère. Un spectre de RMN du
${}^{13}$C trancherait~: deux signaux de carbonyle dans le rapport 2:1
pour mer, un seul signal pour fac.
\textbf{23.} \textbf{L'isomère mer a trois élongations CO \omterm{def:b3:group-theory-applied:activity}{actives en IR},
l'isomère fac deux}~: le produit est l'isomère mer.
\end{solution}
